\section{Framework}
\label{sec:framework}

\subsection{The methylation state space}
\label{sec:statespace}

We represent the epigenetic state of a biological sample as a vector of CpG
methylation beta values. For a sample measured on the Illumina 450K or EPIC array,
this is a point $x \in [0,1]^p$ where $p$ denotes the number of CpG sites passing
quality control, typically $p \approx 400{,}000$--$850{,}000$. We refer to this
space as the \emph{methylation state space} $\mathcal{M} = [0,1]^p$.

Although $\mathcal{M}$ is high-dimensional, empirical methylation data occupies a
much lower-dimensional structure within it. This is the empirical basis for the
success of dimensionality reduction methods in the clock literature and motivates
the manifold assumption we adopt: the biologically relevant states form a smooth,
low-dimensional submanifold of $\mathcal{M}$.

\paragraph{Weighted inner product.}
Raw Euclidean distance in $\mathcal{M}$ treats all CpG sites equally, which is
inappropriate for cross-context comparisons: a change of $0.1$ at a highly variable
CpG site carries different information than the same change at a nearly invariant
site. We therefore equip $\mathcal{M}$ with a weighted inner product. For two
displacement vectors $u, v \in \mathbb{R}^p$, define
\begin{equation}
  \langle u, v \rangle_c = \sum_{j=1}^{p} \frac{u_j v_j}{\sigma_{c,j}^2}
  \label{eq:inner}
\end{equation}
where $\sigma_{c,j}^2$ is the variance of CpG site $j$ across control individuals
in context $c$. This normalises displacements by per-site baseline variability,
making them dimensionless and comparable across tissue types and cohorts.
CpG sites with very low variance (near-constant across individuals) are downweighted
and effectively excluded. Saturation effects near $\beta = 0$ and $\beta = 1$ are
noted as a limitation; a Riemannian extension using the arc-sine or logit
transformation is a natural refinement left for future work.

\subsection{Interventions: general definition}
\label{sec:interventiondef}

\begin{definition}[Intervention]
An \emph{intervention} $I$ is any externally applied perturbation to a biological
sample. Formally, $I$ induces a map $f_I : \mathcal{M} \to \mathcal{M}$, where
$f_I(x)$ denotes the methylation state of a sample that had state $x$ prior to
the application of $I$.
\end{definition}

This definition makes no assumption about the direction or biological meaning of
the perturbation. Interventions include geroprotective candidates (caloric
restriction, rapamycin), known age-accelerating exposures (tobacco smoking,
cytotoxic chemotherapy), pharmacological agents with unknown aging effects, and
environmental or behavioural exposures of any kind. Classification of an
intervention as geroprotective, age-accelerating, or clock-gaming is the
\emph{output} of the framework, not an input to it.

The \emph{mean displacement} induced by intervention $I$ in context $c$ at baseline
state $x$ is
\begin{equation}
  \Delta_{I,c}(x) = \mathbb{E}\bigl[f_I(x) - x \;\big|\; \text{context } c\bigr]
  \label{eq:displacement}
\end{equation}
estimated empirically from paired pre- and post-intervention methylation
measurements in context $c$.

\subsection{The aging direction}
\label{sec:agingdirection}

To identify a clock-independent aging direction, we require a set of interventions
whose biological direction is established by evidence independent of any epigenetic
clock.

\paragraph{Calibration anchors.}
Let $\mathcal{I}^+$ denote a set of interventions with strong prior evidence of
accelerating biological deterioration, and $\mathcal{I}^-$ a set with strong prior
evidence of slowing it, where in both cases the evidence is drawn from mortality,
disease incidence, and healthspan outcomes rather than from clock readings. In the
empirical analysis we use:
\begin{itemize}
  \item $\mathcal{I}^+$: tobacco smoking (epidemiological evidence for lung cancer,
    cardiovascular disease, and all-cause mortality); cytotoxic chemotherapy
    (accelerated functional decline and senescence burden in treated patients).
  \item $\mathcal{I}^-$: caloric restriction (extended healthy lifespan in multiple
    species; the CALERIE randomised trial in humans \cite{belsky2023calerie});
    aerobic exercise (reduced all-cause mortality in observational and intervention
    studies).
\end{itemize}

The labeling sets $\mathcal{I}^+$ and $\mathcal{I}^-$ are calibration anchors, not
a definition of what counts as an intervention. Any perturbation, labeled or
unlabeled, can be classified by the framework once the aging direction is identified.

\paragraph{The optimization.}
The \emph{aging direction} at state $x$ is the unit vector $v^*(x)$ that maximally
separates the displacements of accelerating interventions from those of geroprotective
ones, averaged across contexts:
\begin{equation}
  v^*(x) = \operatorname*{arg\,max}_{|v|_c = 1}
  \left[
    \sum_{I \in \mathcal{I}^+} \sum_c \langle \Delta_{I,c}(x),\, v \rangle_c
    \;-\;
    \sum_{I \in \mathcal{I}^-} \sum_c \langle \Delta_{I,c}(x),\, v \rangle_c
  \right]
  \label{eq:agingdir}
\end{equation}
where $|\cdot|_c$ denotes the norm induced by $\langle \cdot, \cdot \rangle_c$.

\paragraph{SVD formulation.}
Let $A$ be the matrix whose rows are the signed, normalised mean displacements:
for each pair $(I, c)$, the corresponding row of $A$ is $+\Delta_{I,c}$ if
$I \in \mathcal{I}^+$ and $-\Delta_{I,c}$ if $I \in \mathcal{I}^-$, with each
entry divided by $\sigma_{c,j}$. Then $v^*(x)$ is the first right singular vector
of $A$, and the first singular value $\lambda_1$ measures the total signed agreement
across interventions and contexts. The ratio $\lambda_1 / \lambda_2$ serves as a
test of one-dimensionality of the aging signal; we discuss this in
Section~\ref{sec:identification}.

\subsection{The aging trajectory}
\label{sec:trajectory}

\begin{definition}[Aging trajectory]
The \emph{normal aging trajectory} $\gamma : [0,1] \to \mathcal{M}$ is the integral
curve of the aging direction field $v^*$: the unique smooth curve satisfying
$\gamma'(t) \propto v^*(\gamma(t))$ for all $t$, initialised at the mean methylation
state of the youngest individuals in the population and oriented so that $t$
increases in the direction of $v^*$.
\end{definition}

In practice $\gamma$ is estimated as a principal curve \cite{hastie1989} fit to
cross-sectional methylation data from individuals across the age range, initialised
along $v^*$ to ensure correct orientation. The use of chronological age here is
limited to initialisation and orientation of the curve; it plays no role in defining
the direction $v^*$ itself.

\paragraph{Arc-length parameterisation and biological age.}
We parameterise $\gamma$ by arc-length under $\langle \cdot, \cdot \rangle_c$,
denoting the arc-length coordinate as $s \in [0, S]$ where $S$ is the total
length of the estimated trajectory. The \emph{biological age position} of a sample
with methylation state $x$ is
\begin{equation}
  s(x) = \text{arc-length}\bigl(\pi(x)\bigr)
  \label{eq:bioage}
\end{equation}
where $\pi(x) = \operatorname{arg\,min}_{y \in \gamma} \|x - y\|_c$ is the
projection of $x$ onto $\gamma$.

\paragraph{The two-component structure of biological age.}
Biological age as defined above is a position on a curve. For a longitudinally
tracked individual with methylation states $x(t_1)$ and $x(t_2)$ at times $t_1$
and $t_2$, the \emph{biological aging rate} over that interval is
\begin{equation}
  \dot{s} = \frac{s(x(t_2)) - s(x(t_1))}{t_2 - t_1}
  \label{eq:agingrate}
\end{equation}
Position $s$ and rate $\dot{s}$ are independent quantities. An intervention may
affect $s$ (repositioning the cell along the trajectory), $\dot{s}$ (changing the
speed of traversal), or both. Individual variation in natural aging speed is
variation in $\dot{s}$, not variation in $\gamma$: the trajectory is a
population-level geometric object; individual dynamics are parameterised on top
of it.

This decomposition resolves a structural ambiguity in the existing clock literature.
Static clocks (Horvath, GrimAge) are trained to predict quantities correlated with
position $s$. Pace-of-aging measures (DunedinPACE) are trained on longitudinal
biomarker data and therefore estimate $\dot{s}$. These are not competing measures
of the same quantity; they measure genuinely different components of biological age.
We return to this in Section~\ref{sec:discussion} in the context of the CALERIE
trial results.

\subsection{Intervention classification}
\label{sec:classification}

With the aging trajectory $\gamma$ and the projection $\pi$ in hand, we can classify
any intervention formally.

\begin{definition}[Geroprotective intervention]
Intervention $I$ is \emph{geroprotective} at state $x$ if
$s(f_I(x)) < s(x)$: it displaces the sample toward the younger end of the aging
trajectory.
\end{definition}

\begin{definition}[Age-accelerating intervention]
Intervention $I$ is \emph{age-accelerating} at state $x$ if
$s(f_I(x)) > s(x)$: it displaces the sample toward the older end of the aging
trajectory.
\end{definition}

\begin{definition}[Clock-gaming intervention]
Intervention $I$ is \emph{clock-gaming} at state $x$ if the displacement
$f_I(x) - x$ has large magnitude but small arc-length change. Formally, $I$ is
clock-gaming if
\begin{equation}
  \frac{|\pi(f_I(x)) - \pi(x)|_c}{|f_I(x) - x|_c} < \epsilon
  \label{eq:clockgaming}
\end{equation}
for a threshold $\epsilon \ll 1$: the intervention moves the sample substantially
in methylation space but primarily in directions orthogonal to $\gamma$, leaving
biological age position approximately unchanged.
\end{definition}

\begin{definition}[Curve-deforming intervention]
Intervention $I$ is \emph{curve-deforming} if it alters the shape of $\gamma$
itself rather than repositioning a sample along it. Partial cellular reprogramming
via Yamanaka factors is the primary biological example: such interventions do not
move a cell along the normal aging trajectory but place it on a qualitatively
different trajectory. A full treatment of curve-deforming interventions requires
an extended framework and is left for future work; we note the category here for
completeness.
\end{definition}

Note that these categories are not mutually exclusive in general: an intervention
may have both a tangential component (affecting $s$) and an orthogonal component
(clock-gaming). The classification above identifies the dominant character of each
intervention; a quantitative decomposition is given in Section~\ref{sec:results}.

\subsection{Desiderata derived from the construction}
\label{sec:desiderata}

Rather than proposing desiderata and constructing a definition to satisfy them,
we derive desiderata from the construction and show which existing clocks violate
them and why. This is a stronger logical structure: the desiderata are theorems
about the framework, not assumptions.

\paragraph{Desideratum 1: Context invariance.}
A valid measure of biological age should assign the same biological age to a sample
regardless of which tissue or cohort the measurement is taken from.

The aging direction $v^*$ is defined by cross-context agreement: it is the direction
most consistently induced across tissues and cohorts by the calibration interventions.
By construction, the arc-length position $s(x)$ is therefore tissue-independent in
expectation. Clocks trained on blood methylation violate this desideratum when applied
to other tissues, a well-documented failure mode \cite{horvath2013}.

\paragraph{Desideratum 2: Intervention distinguishability.}
A valid measure should distinguish geroprotective interventions from clock-gaming ones.

Under the framework, geroprotective interventions reduce $s(x)$ while clock-gaming
interventions leave $s(x)$ approximately unchanged. These are definitionally distinct.
Existing clocks fail this desideratum: a compound that specifically demethylates
the CpG sites a clock was trained on will reduce the clock reading regardless of
whether $s(x)$ changes.

\paragraph{Desideratum 3: Directional consistency.}
A valid measure should consistently rank known age-accelerating exposures above
baseline and known geroprotective interventions below, across independent datasets.

This desideratum is empirically testable and we report results in
Section~\ref{sec:results}. It is satisfied by construction for the calibration
anchors; its empirical content lies in whether it generalises to novel interventions
not used in the identification step.

\paragraph{Desideratum 4: Monotonicity with health outcomes.}
Position $s$ along $\gamma$ should be monotonically associated with all-cause
mortality in held-out data.

This desideratum connects the formal framework to the clinical literature and is the
primary external validation criterion. It is not satisfied by construction and must
be tested empirically. We report this test in Section~\ref{sec:results}.

\subsection{Relationship to existing clocks}
\label{sec:clockrelationship}

The framework implies a precise criterion for clock validity: a clock $\hat{a}(x)$
is a valid biological age measure to the extent that its output is a monotone
function of arc-length position $s(x)$. Equivalently, the weight vector of a linear
clock should be approximately parallel to $v^*$ in the weighted inner product.

Clock disagreement --- the empirical phenomenon documented in the companion paper
\cite{vinogradova2026desiderata} --- is reinterpreted under this framework as
a measure of how far each clock's weight vector deviates from $v^*$. Clocks whose
weights are more aligned with $v^*$ are better biological age measures; clocks whose
weights have large orthogonal components are measuring tissue-specific, cohort-specific,
or training-artefact variation rather than position on the aging trajectory.

The causality-enriched clocks of Ying et al.\ \cite{ying2024} (DamAge and AdaptAge)
select CpG sites with Mendelian randomization evidence for causal links to aging
traits. Our framework predicts that these sites should be more closely aligned with
$v^*$ than standard clock sites, since causal upstream sites are more likely to
reflect genuine aging-direction variation than correlational ones. We test this
prediction in Section~\ref{sec:results}.

Regarding other molecular substrates: telomere length, proteomic age, and glycan
age are used in our validation as independent evidence for the biological relevance
of $s$, not as inputs to its construction. We note that concordance between
independent imperfect proxies supports but does not validate any individual proxy;
mortality prediction on held-out data is the appropriate validation criterion.
